\section{Autonomous services and data pipeline}\label{sec:7}

\subsection{The inner loop, and why it is deliberately plain}\label{sec:7.1}

\texttt{dcv-controller} runs \textbf{once a second} per room. It reads
\texttt{readings/\{zone\}/co2} and publishes \texttt{commands/\{zone\}/damper}, through a rule with
three parts:

\begin{itemize}
\tightlist
\item \textbf{A hysteresis band.} Above \textbf{800 ppm} the damper follows a proportional ramp;
  below \textbf{650 ppm} it returns to the 0.15 minimum; \emph{between} the two it holds whatever
  position it has. The band is what stops a vane chattering as noise crosses one threshold.
\item \textbf{A ramp, not a switch.} From the minimum at 800 ppm the opening rises linearly to
  fully open at the \textbf{1000 ppm comfort limit}: position
  $= 0.15 + \frac{C - 800}{200}(1 - 0.15)$, clamped open at and above the limit. So a room at the
  limit is already asking for everything the plant can give it.
\item \textbf{A freshness check and a publish dead band.} A reading staler than three sample
  intervals (3 s) produces no command at all: the position is held, nothing is announced. A
  command that differs from the current position by less than \textbf{0.02} is not published,
  because a stream of no-op commands would bury the step the observer is looking for.
\end{itemize}

A room whose \texttt{strategy/\{zone\}} says \texttt{DEGRADED} is not commanded at all, and a room
in \texttt{COMPENSATING} follows the supervisor's trim with CO₂ getting no vote. That is enforced
twice: once by an explicit guard in the control loop, and once by the entity invariant. Asking a
degraded \texttt{VentilationZone} for a normal control command throws.

This is the simplest thing that meets the requirement, and four reasons make simplicity the right
choice rather than a shortcut. First, the outer loop has to \emph{model} the inner one, and an
observer predicting a model-predictive controller inherits that controller's complexity as model
error, which is exactly what the residual cannot distinguish from a fault. Second, and sharpest,
\textbf{integral action would mask the fault}: a PID controller facing a room that will not respond
winds up its integral term and pushes harder, which is precisely how a competent controller hides a
stuck damper from everything downstream. The plain rule commands what the concentration asks for and
nothing more, so the gap between command and effect stays visible. Third, a plain rule is verifiable,
and it is the component that must keep working when everything cleverer has crashed. Fourth, demand
control is a solved problem and improving it would earn nothing the course grades.

Two thresholds are deliberately \emph{not} configurable: the 1000 ppm comfort limit and the 0.15
damper minimum come from the shared topology rather than from deployment config, because letting a
deployment override either would let configuration move the goalposts the evaluation measures.

\subsection{The observer: judging what the command owed}\label{sec:7.2}

The autonomous service proper is \texttt{fault-monitor}, and its core is the \textbf{observer}.
Each damper command is a step. The observer integrates the same mass balance of Section \ref{sec:6.3},
driven by the \emph{commanded} position and the air handler's state, and predicts the CO₂
trajectory that step should have produced over roughly the following ten minutes. The
\textbf{residual} is $r = C_\text{measured} - C_\text{predicted}$: positive means the room is
dirtier than the command should have left it.

\textbf{The observer uses no occupancy, and that is the answer to the project's central risk.} A
prediction needs a generation term, and head-count is the one thing a real building does not
measure. Three design choices together replace it:

\begin{enumerate}
\tightlist
\item \textbf{Judge only commanded step responses.} The \emph{shape} of a step response depends on
  airflow and volume, not on how many people are in the room; a constant occupancy shifts the
  level, not the response to actuation. Judgement is restricted to a \textbf{10 min window} after a
  command, and only a command that moves the damper by at least \textbf{0.05} arms that window. A
  nudge smaller than five points of travel is not a step response worth judging.
\item \textbf{Estimate generation from quasi-steady state, then stop updating it.} With
  $dC/dt = 0$ the balance gives $G_\text{est} = Q\,(C_\text{measured} - C_\text{out})$ directly
  from quantities the system already has. It is smoothed by a first-order filter at gain
  \textbf{0.3} and is updated \textbf{only while the observer is not judging}: not inside an open
  window, not in \texttt{SUSPECT}, not while a case is latched. So occupancy learned before a step
  informs the prediction and occupancy during the step cannot contaminate it.
\item \textbf{Guard with a design-occupancy ceiling.} The estimate is invalid exactly when
  quasi-steadiness is: a room filling up mid-window produces a residual the damper did not cause.
  The guard is a \textbf{second trajectory integrated in parallel} at the room's design occupancy,
  which the observer recovers from the damper authority alone as
  $n_\text{design} = (A_i - 0.35\,\text{area})/7$, inverting the ventilation guideline of Section
  \ref{sec:6.1}. Suspicion is measured against \emph{that} ceiling, so an unmodelled arrival cannot by itself
  confirm a fault.
\end{enumerate}

Between commands the prediction is pulled toward the measurement by a first-order filter at gain
0.3 rather than snapped onto it. That choice is worth a sentence because it is quantitative: a hard
re-anchor would inject a full $\sigma$ of sensor noise into the prediction at every sample, which
raises the residual's standard deviation by $\sqrt{2}$ and, by the run-length formula of Section
\ref{sec:7.3}, would turn a 79-day in-control run length into roughly half a day. The same reasoning decides
what a step anchors to: the \emph{filtered} prediction, not the single raw sample taken at the
moment of the command.

That guard is not free, and Section \ref{sec:13.7} states its cost: in A145 (20 seats) or A115 (10 seats), a
room holding two people has a ceiling far above anything a stuck damper produces, so the fault is
invisible until the room fills, and detection delay therefore depends on occupancy. The harness
prices that dependency at its two ends rather than describing it: with every room at design occupancy
the p95 confirmation delay is \textbf{6 min}, and on a drawn weekday (the same fault, the same
limit, the same seeds) it is \textbf{257 min} (§\ref{sec:11.2}). The closed-loop arm turns out to carry the
points between them, because its 20 runs are 20 \emph{different} drawn weekdays and the harness
records the head count present at each first ignored command. \textbf{Appendix \ref{sec:F.2}} gives that
stratification and states plainly what kind of evidence it is: a stratification of the runs that
happened, not a sweep of a controlled input.

\subsection{Residual, CUSUM, state machine, classifier}\label{sec:7.3}

\textbf{The statistic} is a \textbf{two-sided CUSUM} on the residual. Its upper arm, which every
figure in this report is scored on, is:

\[ S_0 = 0, \qquad S_t = \max\!\big(0,\; S_{t-1} + r_t - k\big), \qquad \text{alarm when } S_t \ge h, \]

where $r_t$ is the residual at cycle $t$, and \emph{which} residual is the subtlety below, since
it is not the published one, with $k = 30$ ppm $= 1\sigma$ the reference value, $h = 150$ ppm
$= 5\sigma$ the limit, and $\sigma$ the sensor noise of Section \ref{sec:6}. The reference value $k$ is set to
$\sigma$ because a stuck damper produces a large growing bias rather than a subtle shift, which
makes the statistic insensitive to noise and to small model errors; subtracting $k$ at every sample
is also what makes a noise tail draw decay away instead of accumulating. The limit $h$ comes from
R-3 through Siegmund's approximation, which the monitor carries in code so the tuning can be
re-derived rather than remembered: at $(k,h) = (1\sigma, 5\sigma)$ the in-control average run length
is $\approx 1.1\times10^{5}$ samples $\approx 79$ days at 60 s, an order of magnitude inside R-3's
one false alarm per week. The derivation came first and the spike's 20-seed check set then
confirmed it, 0 of 20 at 5σ against 1 of 20 at 4σ (§\ref{sec:11.1}); the empirical result is
corroboration, not the reason for the value. Why the CUSUM rather than a raw threshold:
accumulation catches a
\emph{small persistent} bias long before any single-sample threshold fires, which is the regime a
partly stuck damper lives in.

\textbf{The lower arm} is the mirror, $S^{-}_t = \max\!\big(0,\; S^{-}_{t-1} - r_t - k\big)$, and it is
\emph{guarded}: it accumulates against the 420 ppm outdoor floor rather than against the prediction,
so it sees a sensor drifting downward, which under-ventilates an occupied room, without
accusing a room that has simply emptied. Section \ref{sec:13.1} states what that guard can and
cannot catch, and Appendix \ref{sec:F} accounts for the nine of twenty-six pre-fix false confirmations that
came from this arm before ADR-0010 corrected the plausibility guards.

\textbf{One subtlety that matters, and that I would not have guessed before building it: the
residual that is published is not the residual that is accumulated.} The residual on
\texttt{residuals/\{zone\}} is the honest one, measured minus predicted, because that is what a
human and the dashboard should see. The CUSUM is driven instead by measured minus the
\emph{plausibility ceiling} of Section \ref{sec:7.2}, which is the conservative quantity: it is what makes an
unmodelled arrival unable to confirm a fault. The ceiling residual is \emph{subtracted}, not clipped
at zero, because clipping would rectify the sensor noise, give the statistic a positive mean and
destroy the in-control run length the previous paragraph just derived. And that run length is the
entire basis of R-3.

\textbf{Two thresholds, not one.} Opening a case uses a lower bar than confirming it: the statistic
passing \textbf{90 ppm} (0.6 $h$) moves a room to \texttt{SUSPECT}, and only the full $h$ held for
\textbf{two consecutive cycles} confirms. Persistence is the real defence against a tail draw: a
single excursion is gone on the next sample, since the CUSUM subtracts $k$ from every one, while a
stuck damper clears $h$ every cycle. A suspect room whose statistic has decayed back to
\textbf{zero for three consecutive cycles} returns to \texttt{HEALTHY}. Rooms need
\textbf{three usable cycles} to leave \texttt{WARMUP}, and
readings staler than \textbf{three sample intervals} (180 s at the 60 s monitor cadence) drive
\texttt{BLIND}.

\textbf{The state machine} runs one instance per room:
\texttt{WARMUP → HEALTHY → SUSPECT → CONFIRMED → COMPENSATING → SERVICED → WARMUP}, with
\texttt{BLIND} reachable from \texttt{WARMUP}, \texttt{HEALTHY} and \texttt{SUSPECT} on readings
stale beyond three intervals. Four properties of its shape are the design:

\begin{itemize}
\tightlist
\item \textbf{\texttt{BLIND} is not a fault.} It reports unavailable and issues no diagnosis (R-4),
  and there is \emph{no edge} from it into \texttt{SUSPECT} or \texttt{CONFIRMED}: missing data is
  not evidence of a fault. It always returns to \texttt{WARMUP}, never into an open case, so the step
  response is re-identified before judging resumes.
\item \textbf{\texttt{SUSPECT} can recover; \texttt{CONFIRMED} cannot.} \texttt{SUSPECT → HEALTHY}
  is the machine's one recovery edge, taken when $S$ has decayed back to zero for three
  consecutive cycles. Falling below $h$ is not enough: that only stops the confirmation count,
  and a statistic parked between zero and $h$ keeps the case open. The asymmetry is deliberate.
\item \textbf{\texttt{CONFIRMED} and \texttt{COMPENSATING} are latched}, and the only exit is
  \texttt{serviced()}. That is the masking defence: compensation shrinks the residual that found the
  fault, so a monitor that believed its own shrinking residual would be talked out of a true
  diagnosis by the evidence its own response destroyed.
\item \textbf{The active probe re-confirms but cannot clear.} Every 30 min a latched room's monitor
  \emph{asks} for a damper step (as an \textbf{event}, never a command, which is §\ref{sec:4.3}'s authority
  rule holding even here), and the supervisor issues it, rate-limited to one per room per 15 min and
  per-room because a fan-speed probe would perturb all six. A damper being compensated answers a step
  exactly as a repaired one does, so the outcome is recorded either way and the case stays latched
  until service.
\end{itemize}

Every one of those numbers is a named field with a documented default rather than a literal buried
in a branch, which is what let the register of §\ref{sec:11.6} name what each experiment must re-measure;
\textbf{Appendix \ref{sec:D.1}} names eighteen of \texttt{MonitorParameters}' 25 settings. A nineteenth deserves an
argument rather than a line in a list: the plausibility ceiling runs at \textbf{1.25} times design occupancy, and 25 \% of headroom
above design is a deliberately desensitising choice.

\textbf{Appendix \ref{sec:D.1}} prices that choice rather than leaving it asserted.
\texttt{MonitorParameters.\allowbreak withOccupancyCeiling} lets T-24 sweep the headroom against a control arm
with nothing injected, and the sweep says what tightening the ceiling to 1.00 would buy and cost: one
more closed-loop detection and some speed, paid for with 2 false alarms in 20 fault-free runs where
the shipped 1.25 latches 0 of 20. The headroom is therefore not a free lever between R-1 and R-3, and
below 1.00 it stops bounding occupancy at all. That control arm is also what makes the shipped
setting's \textbf{15 of 20} closed-loop detections real detections rather than R-3's false alarms
landing on a room that happened to have a fault in it.

\textbf{The classifier} is a short decision tree over four signals, with an explicit
\texttt{UNCLASSIFIED} branch that alerts a human rather than acting:

\begin{longtable}[]{@{}
  >{\raggedright\arraybackslash}p{(\columnwidth - 4\tabcolsep) * \real{0.26}}
  >{\raggedright\arraybackslash}p{(\columnwidth - 4\tabcolsep) * \real{0.30}}
  >{\raggedright\arraybackslash}p{(\columnwidth - 4\tabcolsep) * \real{0.44}}@{}}
\toprule
\textbf{Signal} & \textbf{Test} & \textbf{What it rules in or out} \\
\midrule
\endhead
Common-mode count & How many of the six residuals are raised; 5 of 6 is the threshold & Many
together ⇒ plant-wide, so a blocked filter; one alone ⇒ local \\
Plant fan power & Flat within 5 \%, or raised beyond 10 \% of the affinity-law prediction & Power
up per unit of flow ⇒ blocked filter; flat ⇒ the fan is doing what it was told \\
Step amplitude & Against $2\sigma = 60$ ppm & Distinguishes a real response from noise; a drifting
sensor still shows the right \emph{shape} \\
Overnight floor & Within 100 ppm of the 420 ppm outdoor level over 24 h & A wrong floor with an
empty building ⇒ sensor drift, since the room must decay to outdoor \\
\bottomrule
\end{longtable}

The tree itself is four lines, evaluated in order, where \emph{localised} means simply ``not common
mode'':

\begin{footnotesize}
\begin{verbatim}
common mode AND power up                              -> BLOCKED_FILTER
localised AND step shape unchanged AND floor wrong     -> SENSOR_DRIFT
localised AND power flat AND step shape NOT unchanged  -> STUCK_DAMPER
otherwise                                              -> UNCLASSIFIED
\end{verbatim}
\end{footnotesize}

Two properties of this design are deliberate. First, \texttt{UNCLASSIFIED} is a \emph{destination},
not an error path: it confirms the residual and declines to name a cause, producing a work order
with an open question rather than a confident wrong answer. Second, the classifier gathers evidence
\textbf{for and against all three hypotheses regardless of which one wins}, so the work order a
technician reads carries the reasoning rather than the verdict, including, for instance, the
measured plant power excess even when no signature matched. Localisation comes from \emph{comparing
rooms}, not from a fourth invented signal: that was an explicit instruction in the risk register,
and Section \ref{sec:13.1} reports what following it cost, which turns out to be more than I expected.

\textbf{What the tree actually emits, measured (T-15).} The table above is what the tree is
\emph{for}; \texttt{ConfusionMatrixExperiment} measures what it does, over \textbf{218 runs} of the
four injectable faults plus a no-fault arm. The headline is the least flattering result in this report
and belongs in the body: \textbf{one of the four mechanisms is classified correctly}. Over seeds 42
and 1--10, \texttt{STUCK\_DAMPER} is named correctly in 8 of 11 runs; \texttt{BLOCKED\_FILTER} comes
out \texttt{UNCLASSIFIED} 11 of 11, exactly as Section \ref{sec:4}'s cross-room analysis predicts, with its
fan-power excess still attached so the work order points at the plant; \texttt{SENSOR\_DRIFT} is
\textbf{never emitted: 0 of 76 runs from a ramp and 0 of 76 from a step offset}, because its leaf
wants an informative overnight floor \emph{and} a step response that still matched in the same
instant and the two injectable sensor mechanisms supply one each and never both; and the stuck-damper
leaf fires \emph{with no step on record}, because a probe is requested through the evidence on
\texttt{FaultStateChanged}, so no step the monitor \emph{asked for} has ever been judged when it
classifies; whether an ordinary control command happened to arm a window by then is luck. That
last one is deliberate rather than an oversight (requiring positive evidence of an unanswered step
was tried and measured, and it left the tree classifying nothing at all), so every diagnosis reached
without a judged step now says in its evidence that it rests on two signals rather than three and that
a probe has been requested. Reordering the pipeline so the probe lands first was built and measured as
T-35 and is \textbf{not adopted}: the third signal is blind after a confirmation rather than merely
late (§\ref{sec:13.6}). \textbf{Appendix \ref{sec:D.2}} carries the matrix, the three off-diagonal mechanisms with the
arithmetic that closes each, and T-35's 176 paired runs.

\subsection{Compensation: the trim is the gain}\label{sec:7.4}

This is the strongest technical result in the project, and it is a closed form rather than a
search. The question is: a damper is frozen partly open, so that room is starved. The only levers
are the other five dampers and one shared fan speed. Can the faulted room be given more air
\emph{without} taking any from the five healthy rooms?

Write $a_i = p_i A_i$ for each opening, $f$ for the frozen room's opening, $h$ for the sum of the
healthy openings at the baseline and $A_0 = h + f$ for the total, then trim \textbf{every healthy
damper by one common factor} $t \in (0,1]$ and raise the fan from $s_0$ to $s$. \textbf{Appendix \ref{sec:E.1}}
carries the four lines of algebra; two facts fall straight out of them and together they are the
whole mechanism.

\textbf{First: one fan speed holds every healthy room's flow at once.} The healthy flow ratio
contains no room index, because the trim is \emph{common} and the split is \emph{proportional}, so
every healthy room's share of the total opening changes by the same factor and a single scalar
$s = s_0\,A(t)/(t\,A_0)$ holds all five at their baseline simultaneously. There is nothing per-room
to balance and no iteration, which is why the supervisor publishes one fan-speed command and five
identical trims rather than solving a five-dimensional problem.

\textbf{Second: at that speed the frozen room's flow is multiplied by exactly $1/t$.}
\textbf{The trim \emph{is} the gain.} Trim the healthy dampers to a quarter and the starved room gets
four times the air, with the other five rooms' airflow bit-for-bit unchanged.

\textbf{Third, and the reason this is deployable: $Q_\text{nom}$ and $\beta$ cancel.} Both appear
in $Q_i$ and in $Q_i$'s baseline, so they vanish from both ratios. The supervisor therefore
never needs to observe the filter blockage or know the plant's nominal flow to compensate
correctly. It needs only the openings and the baseline speed, which it already commands. A
compensation scheme that required $\beta$ would require diagnosing a second fault before answering
the first.

\textbf{Where it runs out.} Two bounds on $t$, both physical rather than numerical, and one case
no lever reaches at all. All three reported with R-2 rather than buried:

\begin{itemize}
\tightlist
\item \textbf{The dampers' floor.} No healthy damper may go below its own minimum, so
  $t \ge \max_i(\text{minimum}_i / p_i)$. With the healthy dampers modulating around 0.60 and the
  floor at 0.15 the trim bottoms out at 0.25, which caps the gain near fourfold; the solver's
  \texttt{MAXIMUM\_REQUESTED\_GAIN} is 4.0 for that reason and not as a tuning knob.
\item \textbf{The fan's ceiling.} $s(t) > s_0$ always, so holding the healthy flows is paid for out
  of fan headroom, and $s$ is capped at 1.0. The plant runs normally at $s_0 = 0.7$, which leaves
  43 \% more air to call on, and that is what makes the lever a lever. When the plant is
  \emph{already} at full speed there is no headroom at all, and the solver neither abandons
  compensation nor quietly starves the healthy rooms: it trims only as far as
  \texttt{HEALTHY\_FLOW\_TOLERANCE} = 0.05 permits, so each healthy room gives up at most 5 \% of
  its airflow, and it \emph{reports the smaller gain it actually achieved}. That 5 \% is kept on a
measurement rather than assumed: the over-supply side was swept from 0.05 to unbounded and bought
R-2 nothing (ADR-0009, §\ref{sec:11.2}). Flow is then
  redistributed rather than created, which is the honest description of a saturated fan.
\item \textbf{A damper stuck fully shut is unreachable.} The gain multiplies an airflow of zero, so
  no trim and no fan speed touches it. The solver returns \texttt{CompensationPlan.none} saying so,
  and the supervisor raises a \emph{critical} work order with occupancy advice rather than
  pretending. This is R-2's stated bound, and Section \ref{sec:13} concedes the uncomfortable part: ADR-0008
  defines ``stuck shut'' \emph{as} stuck at 0.15, so R-2's scenario is pinned at the one stuck-shut
  position that is reachable. The unreachable case is excluded by a parameter, not by physics.
\end{itemize}

\textbf{The second lever, and why the trim alone was not enough.} The derivation above assumes there
is a trim to make, and in the regime R-2's drawn weekday covers there is not: the healthy rooms are
below the inner loop's 800 ppm opening threshold when a plan is solved, so they are already on their
0.15 floor, the trim floor is $t = 1$, and the closed form returns the plant unchanged, a
supervisory loop worth nothing and, at design occupancy, less than nothing (§\ref{sec:11.2}). The split law
gives a second lever, because $Q_i$ is linear in $s$, so raising the fan multiplies \emph{every}
room's flow including a frozen room's. But that lever is \textbf{not selective}: with no trim in
force the frozen room's gain and the healthy rooms' over-supply are the same number, so ADR-0008's
43 \% of headroom above $s_0 = 0.7$ is spendable selectively only \emph{through} a trim, and without
one the solver may spend it only as far as \texttt{HEALTHY\_FLOW\_TOLERANCE} = 0.05 permits.
\textbf{Appendix \ref{sec:E.2}} derives the bound and prices the raise, which is not free: fan power goes as
$s^3$, so 0.70 to 1.00 costs 2.92 times the fan power, and the work order states that multiple
rather than implying it.

\textbf{The measured instance.} At the parameters of Section \ref{sec:6}, with the five healthy dampers at
0.60, the fan at its normal 0.70 and A145's damper frozen at the 0.15 floor, A145 receives
\textbf{59.69 l/s} against its 165.9 l/s design flow. The solver trims the healthy dampers to
$t = 0.360$ (0.60 → 0.216) and raises the fan to $s = 0.880$, which restores A145 to
\textbf{165.9 l/s, its full design flow}, while the five healthy rooms' flows are unchanged
\textbf{to within $10^{-15}$}: eight decimal places and more.

\textbf{Neither bound binds here, and that is the air handler's doing rather than the solver's.}
An earlier revision of this report read the same case off the \emph{superseded} 0.42 m³/s plant and
found the trim bottoming out at its floor just as the fan reached its ceiling: an elegant
coincidence, and a property it attributed to the topology. On the plant the repository now ships,
0.59 m³/s, the two bounds are not close: the trim stops at 0.360 against a floor of \textbf{0.25}
(the tightest healthy damper's $0.15/0.60$), and the fan at 0.880 against a ceiling of
\textbf{1.00}. What changed is not the closed form but the plant it is evaluated on. Re-sizing the
air handler to cover its design condition at the speed it actually runs (Section \ref{sec:13.3}) bought
compensation margin that the earlier sizing did not have, and the coincidence was an artefact of
being short of air rather than a property of the topology.

\textbf{That margin is this case's and must not be read as the system's.} This is the single-room
frozen-damper case with five healthy neighbours modulating at 0.60 and the rest of the building
behaving. Section \ref{sec:11.2}'s measured multi-room runs at design occupancy are the regime where it runs
out: there the healthy dampers are already near their own floors when a plan is solved, the fan
reaches \textbf{1.00} on one of five seeds with every zone faulted and no damper left to trim, and
the achieved gains run \textbf{1.05 to 1.77} against the 2.3 the comfort rule asks for. Headroom in
the one-room case is not headroom in the regime R-2 fails in. Section \ref{sec:11.3} gives the per-room
numbers; Section \ref{sec:10} names the test that pins them.

\subsection{The data pipeline}\label{sec:7.5}

Readings travel through one path that serves live control, diagnosis, the dashboard and the
evaluation. \texttt{co2-sensor} publishes to \textbf{mosquitto}, which decouples one producer from
four consumers and turns a dead publisher into silence rather than a blocked caller, silence being
exactly what R-4 is written to detect. \texttt{dcv-controller}, \texttt{fault-monitor} and
\texttt{supervisor} each append to \textbf{TimescaleDB}: commands, readings, residuals and
decisions, strategies. The dashboard queries windows from the same store and subscribes to the
broker for live state.

\textbf{Data quality} is judged at the edge that can see it, and there is more of it than the
requirements ask for. \texttt{co2-sensor} stamps every reading \texttt{GOOD}, \texttt{STALE},
\texttt{MISSING} or \texttt{DUPLICATE}: a value identical to the last one is \texttt{DUPLICATE}
until it has been unchanged past the staleness window, when it becomes \texttt{STALE}; a failed read
publishes \texttt{MISSING} \emph{carrying the last known value}, never a fabricated one, and nothing
at all is published before the first successful read. The comparison is exact rather than
tolerance-based, deliberately, because BuildSim round-trips values as fixed-decimal strings and a
tolerance would classify slow genuine drift, a fault, as a duplicate. Downstream, the controller
drops anything not \texttt{GOOD} and drops out-of-order redelivery, and the dashboard's projection
accepts a reading only if it is strictly newer than the one it holds, so QoS-1 redelivery cannot
rewrite history.

The honest gap is not in the mechanism but in the \emph{specification}: no requirement says what
\texttt{quality} must mean, what must happen to a duplicate, or how out-of-order arrivals must be
handled, so the behaviour above is a reasonable invention rather than something a test can hold to
an agreed criterion. That is Section \ref{sec:3.2}'s first gap.

\textbf{What is retained}: everything, for the length of a run, in the hot store only. There is no
cold tier. The cold path the plan scheduled as Parquet (T-22) exists as CSV plus a one-command
conversion instead: the harness writes nine tables totalling \textbf{430\,797 bytes} into
\texttt{evaluation/results/} and generates \texttt{export.sql} beside them, whose nine \texttt{COPY}
statements convert every one of those tables to Parquet, one statement each. Parquet itself is not vendored, and the reason is
measured rather than asserted: \texttt{parquet-hadoop} 1.15.2 with \texttt{hadoop-common} 3.4.1
closes over 98 jars and 64.8 MiB, and \texttt{duckdb\_jdbc} 1.5.6.0 over 1 jar and 81.2 MiB, against
roughly 150 kB of measurement data.
